\scanpage{6}

\uuline{Proposition} On a donc des \deleted{\unsure{hom}} inclusions
\[
\mathcal E_0=\mathcal E_1\subset\mathcal E_2\subset\mathcal E_3\subset\cdots
\]
les $\mathcal E_i$ sont des faisceaux loc. libres de rang $i$ pour $i\geq1$. D'ailleurs
\[
\mathcal E_0=\underline O_S
\]

\uline{Proposition 2} $\mathcal E_0=\mathcal E_1=\underline O_S$, d'autre part pour $i\leq j$ l'injection $\mathcal E_i\longrightarrow\mathcal E_j$ est \deleted{\unsure{locale}} universellement injective.
\deleted{Le faisceau} \inserted{Soit} $\mathcal L_n\simeq\mathcal E_{\unsure{n}}/\mathcal E_{n-1}$ \inserted{(pour $n\geq\unsure{0|2}$, \deleted{\unsure{$\mathcal L_0=\underline O_S$}}) \unsure{les $\mathcal L_i$, soit $i\neq1$}; alors on a des isomorphismes} \deleted{est isomorphe à $\mathcal L^{\otimes\unsure{k}}$, et $\mathcal E_{n+1}/\mathcal E_n$ est isomorphe à $\mathcal L^{\otimes n}$.}
\[
\mathcal L_n\otimes\mathcal L_m\simeq\mathcal L_{m+n}\qquad\text{si }n,m\begin{cases}\geq0\\\neq1\end{cases}
\]
\inserted{provenant de} \deleted{\unsure{faisant de}} \inserted{une structure} sur $\coprod_{n\geq0}\mathcal L_n$ \deleted{\unsure{une}} d'Algèbre graduée comm.

Pour la première assertion, \unsure{avec|grâce à} la platitude coh. on est ramené au cas d'un corps, où c'est trivial car on a une injection.

Notons que $\mathcal E=\varinjlim\mathcal E_n=f_*(\varinjlim\underline L^{\otimes n})$ est une $\underline O_S$-\uline{Algèbre}, dont la connaissance d'ailleurs implique la connaissance de $\underline\Gamma_*^{\unsure{\mathcal E}}(f,\underline L)=\coprod\mathcal E_n$, car $\mathcal E_m\times\mathcal E_n\unsure{\longrightarrow|-}\mathcal E_{m+n}$ est précisément induit par la multiplication de $\mathcal E$. On a, en passant au \unsure{gr. ann.} associé à la filtration croissante donnée, une Algèbre $\coprod_{n\geq0}\underline{\mathcal L}_n$, où les $\underline{\mathcal L}_n\simeq\mathcal E_n/\mathcal E_{n-1}$ pour $n\geq0$, en modifiant pour $n=0$ par $\underline O_S$. Il faut prouver que
\[
\underline{\mathcal L}_n\otimes\underline{\mathcal L}_m\unsure{\longrightarrow|-}\underline{\mathcal L}_{n+m},\qquad n,m\begin{cases}>0\\\neq1\end{cases},\quad\deleted{n,m}
\]
est bien isom., ou encore surjectif, \deleted{\unsure{et pour}} Si $n$ ou $m=0$, \deleted{\unsure{ramenons nous} \illegible{3} \unsure{au voisinage}} c'est trivial, sinon on a $n,m\geq2$, et on voit.

\mspage{5}
